% Source: source/普通高考/2007/2007山东文.pdf, page 1, question 10.
% Node -> directed-edge list: 开始 -> 输入 n -> S=0,T=0 -> n<2?;
% 是 -> 输出 S,T -> 结束; 否 -> S=S+n -> n=n-1 -> T=T+n -> n=n-1
% -> left loop-back to the shared stem before n<2?. All connectors are solid;
% the loop and the output branch are independent, with no extra merge edges.

\begin{tikzpicture}[
    >=Stealth,
    x=1cm,
    y=0.86cm,
    line width=0.45pt,
    line cap=round,
    line join=round,
    font=\normalsize,
    flowtext/.style={anchor=center,align=center,
        execute at begin node={\setlength{\parindent}{0pt}}},
    every node/.style={flowtext},
    flowbox/.style={draw,flowtext,minimum height=0.68cm,
        inner xsep=4pt,inner ysep=2.5pt},
    terminal/.style={flowbox,rounded corners=0.18cm,minimum width=1.35cm},
    process/.style={flowbox},
    decision/.style={flowbox,diamond,aspect=2.9,minimum height=0.95cm,
        minimum width=3.45cm,inner sep=0pt},
    io/.style={flowbox,trapezium,trapezium left angle=72,
        trapezium right angle=108,minimum height=0.78cm},
    branchlabel/.style={flowtext,inner sep=0pt},
]
    \path[use as bounding box] (-4.50,-1.00) rectangle (5.35,10.60);

    \node[terminal] (start) at (0,10.05) {开始};
    \node[io,minimum width=1.85cm] (input) at (0,8.75) {输入 $n$};
    \node[process,minimum width=3.35cm] (init) at (0,7.35)
        {$S=0,T=0$};
    \node[decision] (test) at (0,5.65) {$n<2?$};
    \node[process,minimum width=2.70cm] (sum) at (0,4.20)
        {$S=S+n$};
    \node[process,minimum width=2.65cm] (dec1) at (0,2.90)
        {$n=n-1$};
    \node[process,minimum width=2.65cm] (tadd) at (0,1.60)
        {$T=T+n$};
    \node[process,minimum width=2.65cm] (dec2) at (0,0.30)
        {$n=n-1$};
    \node[io,minimum width=2.55cm] (output) at (3.55,4.05)
        {输出 $S,T$};
    \node[terminal] (finish) at (3.55,2.70) {结束};

    \draw[->] (start.south) -- (input.north);
    \draw[->] (input.south) -- (init.north);
    % Keep the decision-entry arrowhead at least 3 mm from the loop junction.
    \coordinate (merge) at (0,6.60);
    \draw (init.south) -- (merge);
    \draw[->] (merge) -- (test.north);

    % 否 branch updates the running sums and decrements n twice.
    \draw[->] (test.south) -- node[branchlabel,right=3pt] {否} (sum.north);
    \draw[->] (sum.south) -- (dec1.north);
    \draw[->] (dec1.south) -- (tadd.north);
    \draw[->] (tadd.south) -- (dec2.north);

    % 是 branch exits to output and termination.
    \draw[->] (test.east) -- node[branchlabel,above=3pt] {是}
        (3.55,5.65) -- (output.north);
    \draw[->] (output.south) -- (finish.north);

    % Loop-back from the final decrement to the shared stem.
    \coordinate (loop-bottom) at (-3.55,-0.04);
    \coordinate (loop-tip) at (-0.35,6.60);
    \draw[->] (dec2.south) -- (loop-bottom) -- (-3.55,6.60) -- (loop-tip);
    \draw (loop-tip) -- (merge);
\end{tikzpicture}
